\subsection{The Role of Play in Moral Learning and Socialization}
\label{sec:5.7.1}
Children build much of their moral sense through play. Instruction supplies far less of it than adults assume. A child who tests a rule against a playmate — pushing it, breaking it, watching the other child's face — gets an answer no lecture supplies: acceptance or rejection, delivered immediately by someone with no stake in teaching a lesson. Developmental psychologists point to play, more than direct teaching, as the route by which perspective-taking, empathy, cooperation, and respect for rules actually get built, through repeated low-stakes trial and error \autocite{gray2013free}.

Section 4.3.1 covers a related case from the exploration side of play: agents in OpenAI's Hide and Seek environment developed tool use and counter-strategies nobody had programmed, purely through repeated competitive play against each other. That result is about competence emerging from play. What this section asks is narrower and different in kind: not whether play can teach a system to do something, but whether it can teach a system standards for how to treat the agents it plays with.

An AI system has no biological drives. And its witnessing of its "own" winning or losing is not necessarily a felt personal experience. But a training context can still be built to have the structural features that make human-like play work as a moral teacher: actions that carry real consequences for other agents in the environment, feedback that comes from the environment itself, with no supervisor's label attached after the fact, and enough latitude in how the system responds that its choices are its "own" and not a script it is executing. That combination is what would make a scenario play. Whether emotion is present has nothing to do with it.

Two existing games illustrate what a moral version of that structure looks like. Quandary, an ethics-focused game built by the Learning Games Network with MIT's Education Arcade, walks a player through a sequence of dilemmas with no dictated right answer and scores how the player reasons, leaving aside which option they pick \autocite{lgn2012quandary} — exactly the design principle a moral-training environment for an AI system would need: evaluate the deliberation itself. Hanabi, a cooperative card game where each player can see everyone's hand but their own and must communicate through the game's own restricted signals, has become a standard AI research benchmark for the capacities morality runs on directly: modeling what a partner does and does not know, and communicating honestly within a shared, rule-bound system \autocite{bard2020hanabi}. Hide and Seek's payoff was adversarial; Hanabi's is intentionally cooperative, and that difference is close to the whole of what separates this section's territory from section 4.3.1's.

\runin{Theoretical Foundations}
Four frameworks from developmental and social psychology give the play-to-morality route something more specific to aim for than a general appeal to "learning through experience."

\begin{itemize}
  \item Constructivism: Jean Piaget treated play as the process by which a child actively builds understanding through interaction with the world, and does not absorb it ready-made \autocite{piaget1962play}. The AI parallel is a system that refines its own model of a situation's stakes by acting on it repeatedly. The alternative only sorts situations against a fixed label.
  \item Social learning: Albert Bandura's account of behavior acquired by observing and imitating others \autocite{bandura1977social} suggests an AI system could pick up norms of restraint or honesty the way a child does — by watching what other agents in its environment do, and what happens to them, and not only from a reward signal attached directly to its own behavior.
  \item Care ethics: Carol Gilligan's account of moral development centers empathy, relationship, and responsiveness to another's situation over rule-following in the abstract \autocite{gilligan1982different}. It belongs here with its origin attached, because that origin is a correction to this chapter's own foundation: Gilligan built the account as a critique of the Kohlbergian stage model section 5.1.1 leans on, having found that its scoring treated relational reasoning as a lower stage (see the box at section 5.1.1). An AI system that plays repeatedly against or alongside other agents has repeated occasions to register how its actions land on them emotionally as well as materially — a feedback loop closer to responsiveness than to rule compliance.
  \item Theory of mind: modeling what another agent believes, wants, and does not yet know is a precondition for most behavior worth calling moral, and it is exactly what a game like Hanabi is built to demand. Training inside games that reward accurate modeling of a partner's knowledge state builds a specific, checkable version of theory of mind, which is more than a general claim that the system "understands other agents."
\end{itemize}
None of these frameworks was built with machines in mind, and the step from a human child's embodied, emotionally loaded play to a training loop is not one-to-one. They specify what play is doing when it works as a moral teacher. Whether the mechanism transfers to a system intact is a separate question, and these findings do not answer it.

\runin{Environments for Ethical Practice}
Knowing why play might work is a different problem from building it into a curriculum, which means being specific about the environment, the feedback, and the pace.

\begin{itemize}
  \item Simulated scenarios with real stakes: an AI system controlling a simulated vehicle in a modeled city, facing pedestrians, cyclists, and other drivers, meets the same dilemma that gives the trolley problem its bite, but embedded in an ongoing interactive scenario, so that the system experiences the downstream consequences of a choice and does not merely name one.
  \item Structured difficulty: a curriculum that starts with low-stakes dilemmas and escalates only once the system handles the easy cases reliably mirrors how a child's own play scales in difficulty as competence grows. A healthcare AI system, for instance, might start with the mechanics of patient confidentiality before being handed obligations that genuinely conflict with each other.
  \item Mentored practice: feedback on a specific attempt at a task like conflict mediation, instead of a rulebook issued once and left alone, is closer to how a coach corrects a player mid-game than to how a policy is handed down in advance.
\end{itemize}
Multi-agent environments are their own case, and a more developed one: section 5.5.1 already covers cooperative and competitive multi-agent training in depth, including reward structures that favor prosocial behavior and agents that learn to share resources or defer to each other under simulated pressure. What that discussion does not settle, and what a play-based framing has to answer and cannot assume, is whether the resulting behavior is moral learning or a policy that merely looks prosocial inside the reward function that produced it.

None of this makes an AI system's play equivalent to a child's. A child's play is charged by stakes that matter to the child — status among peers, the sting of exclusion — where an AI system's "play" is whatever its designer wired the reward function to care about. The frameworks above are worth building toward because they specify what a route into moral competence through practice would need to include. It is not offered because the analogy between a training loop and a playground is more than partial.
